\documentclass[10pt,a4paper]{article}
\usepackage[margin=18mm]{geometry}
\usepackage{enumitem}
\usepackage[hidelinks]{hyperref}
\setlength{\parindent}{0pt}
\begin{document}
{\LARGE Hyeong-Gon Jo}\\[3pt]
Building energy researcher · scientific software developer\\
Building Simulation Lab, Seoul National University\\
\href{mailto:hyeonggon.jo@snu.ac.kr}{hyeonggon.jo@snu.ac.kr} \quad \url{https://github.com/Gonie-Gonie} \quad \url{https://orcid.org/0009-0004-8821-275X}

\section*{Profile}
I study building energy performance, green retrofit evaluation, HVAC control, and simulation methodology. I also develop small, inspectable tools that make engineering research easier to repeat and share.

\section*{Education}
\begin{itemize}[leftmargin=*,nosep]
\item \textbf{Ph.D. in Architecture and Architectural Engineering}, Seoul National University \hfill 2022.03–2026.02\\\emph{Simplification Methods for Dynamic Simulation Tools to Support Design-Stage Decision-Making}
\item \textbf{M.S. in Architecture and Architectural Engineering}, Seoul National University \hfill 2020.03–2022.02\\\emph{Stochastic Daylight Model and Electric Lighting Control}
\item \textbf{B.S. in Architecture and Architectural Engineering}, Seoul National University \hfill 2016.03–2020.02
\end{itemize}

\section*{Selected Research Projects}
\begin{itemize}[leftmargin=*]
\item \textbf{Performance evaluation of green-remodeled buildings} \hfill 2025.06–2025.10\\Korea Real Estate Board; Researcher · simulation and data analysis. Collected operational data and used dynamic simulation to evaluate realized energy savings and performance gaps in completed retrofit projects.
\item \textbf{Prototype energy and carbon simulator for existing buildings} \hfill 2024.07–2025.12\\Korea Authority of Land \& Infrastructure Safety; Simulation architecture and EnergyPlus workflow. Developed a prototype that converts structured, spreadsheet-oriented inputs into repeatable EnergyPlus models and retrofit scenarios.
\item \textbf{Green-remodeling carbon-reduction factor database} \hfill 2022.12–2024.12\\Korea Authority of Land \& Infrastructure Safety; EnergyPlus modeling, large-scale simulation, and decision-tool development. Built archetype-based reduction factors and long-term emissions estimates through a large simulation campaign for aging buildings.
\item \textbf{Integrated optimization simulator for industrial FAB utility systems} \hfill 2022.05–2025.03\\Samsung C\&T; System modeling, estimation, control optimization, and Python simulator. Developed physics-based and data-driven models for air-handling, heat-recovery, and heat-source systems, then connected them in a node-based simulator.
\end{itemize}

\section*{Selected Publications}
\begin{enumerate}[leftmargin=*]
\item Park, C. H., Lee, S. J., Jo, H. G., Park, C. S. (2026). Implications of geometry simplification in dynamic building energy modeling: Comparative analysis of detailed and simplified geometries. Energy and Buildings, 117839.
\item Kim, J. H., Kim, Y. S., Jo, H. G., Urabe, E., Mun, J., Shin, Y., Park, Y., Park, C. S. (2025). Accuracy, generalizability, and extrapolation ability of physics-based, data-driven, and hybrid models for real-life cooling towers. Building and Environment, 112756.
\item Mun, J., Jo, H. G., Kim, J. H., Urabe, E., Mun, J., Park, Y., Park, C. S. (2025). Causality-driven, interpretable surrogate modeling for DOAS preheating control. Energy and Buildings, 116625.
\item Kim, J. H., Kim, Y. S., Jo, H. G., Mun, J., Park, C. S. (2025). Domain-invariant representation learning for generalizable chiller model: a real-world case study. Energy and Buildings, 116168.
\item Ra, S. J., Jo, H. G., Park, C. S. (2025). Real-time model predictive control of energy recovery ventilators for a school building. Science and Technology for the Built Environment, Vol. 31, No. 9, pp. 1155-1166. https://doi.org/10.1080/23744731.2025.2539048
\item Jo, H. G., Park, C. H., Lee, S. J., Mun, H. J., Park, C. S. (2025). Running EnergyPlus Using Spreadsheet Inputs: Focus on Geometry Information. Proceedings of the Fall Conference of the Architectural Institute of Korea, Vol. 45, No. 2, pp. 574-575. October 29-31, 2025, Science and Technology Convention Center, Seoul.
\item Mun, J., Jo, H. G., Park, C. S. (2025). Toward scalable prediction of indoor thermal dynamics: Neural-network-implemented state-space (NNiSS) model. Energy and Buildings 331, 115359.
\item Jo, H. G., Kwon, S. O., Kim, U. H., Yoo, Y. S., Park, C. S. (2024). Establishment and Application of a Greenhouse Gas Reduction Factor Database for Policy Decision-Making in Green Remodeling Projects. Proceedings of the Fall Conference of the Architectural Institute of Korea, Vol. 44, No. 2, pp. 501-502. October 23-25, 2024, The-K Hotel Gyeongju, Gyeongju, Gyeongsangbuk-do.
\item Jo, H. G., Kim, Y. S., Kim, J. H., Park, C. S., Urabe, E., Mun, J., Shin, Y., Park, Y. (2024). Optimal Control Strategies of a Heat Recovery System for a Building. Proceedings of SASBE 2024, November 7-9, 2024, Auckland, New Zealand. https://doi.org/10.1007/978-981-96-4051-5
\item Jo, H. G., Lee, S. J., Park, C. H., Mun, H. J., Park, C. S. (2024). Running EnergyPlus Using Spreadsheet Inputs: Focus on HVAC Systems. Proceedings of the Fall Conference of the Architectural Institute of Korea, Vol. 45, No. 2, pp. 576-577. October 29-31, 2025, Science and Technology Convention Center, Seoul.
\item Lee, S. J., Jo, H. G., Yoo, Y. S., Park, C. H., Park, C. S. (2023). Issues in Quasi-Steady-State Building Energy Analysis Tools: Focusing on ECO2, the Korean Building Energy Efficiency Rating Certification Evaluation Program. Journal of the Architectural Institute of Korea, Vol. 40, No. 1, pp. 169-176.
\item Jo, H. G., Choi, S. H., Park, C. S. (2023). Linear regression based indoor daylight illuminance estimation with simple measurements based on daylight coefficient method. Journal of Building Performance Simulation 16(5), pp. 574-587. https://doi.org/10.1080/19401493.2023.2185684
\item Jo, H. G., Choi, S. H., Heo, S. Y., Park, C. S. (2022). Calibration Method for a Daylight Simulation Model Using Field Measurement Data. Proceedings of the Fall Conference of the Korean Institute of Architectural Sustainable Environment and Building Systems, pp. 119-121. November 25, 2022, Yonsei University, Seoul.
\item Jo, H. G., Choi, S. H., Park, C. S. (2022). Indoor Daylight Estimation Model Based on the Daylight Coefficient Concept. Proceedings of the Spring Conference of the Architectural Institute of Korea, Vol. 42, No. 1, pp. 399-400. April 28-29, 2022, The-K Hotel, Seoul.
\item Jo, H. G., Choi, S. H., Park, C. S. (2022). Real-time indoor daylight illuminance simulation of an existing building using minimal data. Proceedings of the 2022 Winter Simulation Conference, December 11-14, 2022, Singapore, Singapore.
\item Jo, H. G., Kim, Y. S., Park, C. S. (2021). Application of daylighting Gaussian process emulator to lighting control of an existing building. Proceedings of the 2021 Winter Simulation Conference, December 13-17, 2021, Phoenix, Arizona, USA.
\item Jo, H. G., Ra, S. J., Kim, J. H., Park, C. S. (2021). Real-Time Electric Lighting Control Considering the Uncertainty of a Daylight Model. Proceedings of the Fall Conference of the Architectural Institute of Korea, Vol. 41, No. 2, pp. 332-333. October 27-30, 2021, Yeosu Expo Convention Center, Yeosu, Jeollanam-do.
\item Jo, H. G., Kim, Y. S., Park, C. S. (2021). Stochastic Daylight Model and Its Application to Model Predictive Control in Large-Space Buildings. Journal of the Architectural Institute of Korea, Vol. 37, No. 12, pp. 255-264.
\item Jo, H. G., Park, C. S. (2020). Illuminance prediction using a hybrid simulation model based on a single reference measurement. Proceedings of the 2020 Winter Simulation Conference, December 14-18, 2020, Orlando, USA. Virtual Conference.
\item Jo, H. G., Shin, H. S., Kim, Y. S., Park, C. S. (2020). Prediction of Illuminance Ratios Between Indoor Points Using Gaussian Processes. Proceedings of the Fall Conference of the Korean Institute of Architectural Sustainable Environment and Building Systems, pp. 153-154. November 19-21, 2020, Korea National University of Transportation, Chungju.
\end{enumerate}

\section*{Research Software}
\begin{itemize}[leftmargin=*]
\item \textbf{EPlusSimple} --- A spreadsheet-driven EnergyPlus workflow that validates inputs, builds models, runs simulations, and returns compact result summaries.\\\url{https://github.com/snu-bslab/EPlusSimple}
\item \textbf{semantic-idf} --- A desktop workbench for inspecting, diagnosing, cleaning, comparing, and converting EnergyPlus IDF and epJSON models.\\\url{https://github.com/Gonie-Gonie/semantic-idf}
\item \textbf{oo-docs} --- One structured Python document tree rendered to DOCX, PDF, and HTML with shared data, figures, citations, and layout.\\\url{https://github.com/Gonie-Gonie/oo-docs}
\item \textbf{eng-lang} --- An experimental engineering workflow language that preserves units, schemas, axes, provenance, and review artifacts.\\\url{https://github.com/Gonie-Gonie/eng-lang}
\item \textbf{rusted-energyplus} --- An evidence-driven Rust reimplementation experiment for selected EnergyPlus algorithms and compatibility behavior.\\\url{https://github.com/Gonie-Gonie/rusted-energyplus}
\item \textbf{hvac-studio} --- A component-and-node authoring environment and runtime for interconnected building-system models.\\\url{https://github.com/Gonie-Gonie/hvac-studio}
\end{itemize}

\section*{Selected Awards}
\begin{itemize}[leftmargin=*]
\item 2025-12-02 --- \textbf{Best Paper Award, ISHVAC 2025}, 14th International Symposium on Heating, Ventilating and Air-Conditioning. Causality-driven, interpretable surrogate modelling for DOAS preheating control
\item 2024-10-25 --- \textbf{Best Presentation Paper Award}, Architectural Institute of Korea. Establishment and Application of a Greenhouse Gas Reduction Factor Database for Policy Decision-Making in Green Remodeling Projects
\end{itemize}
\end{document}
